%% kapa: exam=1 round=23 year=2012 period=1 form=A subject=economics subject_ko=경제원론 seq=40 number=80
%% source: Q-Net 감정평가사 기출문제 (cst00309) artlSeq=5052578 file=raw/1008/5052578/2012년 제23회 감정평가사 제1차시험 기출문제(1교시 A형).pdf
%% conversion: kordoc 4.24.51 -> tools/gen1.py (auto); review=proofread-sonnet; flags=none
\begin{question}{80}
합리적 경제주체들이 인플레이션율을 6\%로 예상하고 다음과 같은 경제행위를 하였다. 실제 인플레이션율이 3\%일 때 손해를 보는 경제주체를 모두 고른 것은?
\begin{framedtext}
ㄱ. 고정금리로 정기예금에 가입한 가계 \par
ㄴ. 고정된 봉급의 임금계약을 체결한 근로자 \par
ㄷ. 고정금리로 국채를 발행한 정부 \par
ㄹ. 고정금리로 주택담보 대출을 받은 차입자
\end{framedtext}
\begin{choices}
  \item ㄱ, ㄴ
  \item ㄴ, ㄷ
  \item ㄷ, ㄹ
  \item ㄱ, ㄷ
  \item ㄴ, ㄹ
\end{choices}
\end{question}
\answer{3}
